% Revision of the 27 September reviewed manuscript; additions dated 13 October 2026.
% The new material has not inherited the earlier package's review certification.
\documentclass[11pt]{article}

\usepackage[T1]{fontenc}
\usepackage[utf8]{inputenc}
\usepackage{lmodern}
\usepackage{microtype}
\usepackage{amsmath,amssymb,amsthm,mathtools}
\usepackage{geometry}
\usepackage{booktabs,tabularx,array}
\usepackage{enumitem}
\usepackage{xcolor}
\usepackage[hidelinks]{hyperref}
\usepackage{url}

\geometry{margin=1.05in}
\setlength{\parindent}{1.2em}
\setlength{\parskip}{0.15em}
\setlist{nosep}

\newtheorem{theorem}{Theorem}[section]
\newtheorem{lemma}[theorem]{Lemma}
\newtheorem{proposition}[theorem]{Proposition}
\newtheorem{corollary}[theorem]{Corollary}
\theoremstyle{definition}
\newtheorem{definition}[theorem]{Definition}
\newtheorem{example}[theorem]{Example}
\theoremstyle{remark}
\newtheorem{remark}[theorem]{Remark}

\newcommand{\R}{\mathbb{R}}
\newcommand{\Ftwo}{\mathbb{F}_2}
\newcommand{\bits}{\{0,1\}}
\newcommand{\sign}{\operatorname{sgn}}
\newcommand{\supp}{\operatorname{supp}}
\newcommand{\Ch}{\operatorname{Ch}}
\newcommand{\DC}{\mathcal{E}^{\{M\}}}
\newcommand{\eps}{\varepsilon}
\newcommand{\cA}{\mathcal{A}}

\title{\vspace{-1.2cm}\textbf{Finite-Jet Sign Rigidity on Hyperplane Arrangements\\and Denjoy--Carleman Realization}}
\author{Brian Theory}
\date{13 October 2026}

\begin{document}
\maketitle

\begin{abstract}
Let $\mathcal A$ be a finite real affine hyperplane arrangement. We study which $\{\pm1\}$-valued chamber labelings can occur as the sign of a function with prescribed regularity. A first nonzero homogeneous Taylor term inherits the chamber signs weakly, and factor multiplicities force a constant crossing parity. This gives a short arrangement-level proof of the known smooth signature-flatness obstruction and, in finite regularity, explicit $C^d$ orthant and $C^N$ arrangement rigidity statements.

For normalized compactwise Roumieu Denjoy--Carleman classes $\mathcal E^{\{M\}}$, the Denjoy--Carleman criterion yields an exact realization dichotomy. In the quasianalytic case the realizable chamber laws are precisely the signed parity laws; in the non-quasianalytic case every chamber labeling is realized by explicit selector formulas. Thus an arrangement of $m$ distinct hyperplanes with $r$ chambers admits exactly $2^{m+1}$ and $2^r$ realizable sign laws in the two regimes. We also record corner-margin consequences and a fixed-dimensional polynomial comparison on a cube with coordinate-wall strips of width $\delta$ removed: signed parity has constant exact degree, non-affine laws have degree $\Theta(\log(1/\delta))$, and any nonconstant law requires degree $\Theta(1/\delta)$ under a global amplitude bound and fixed margin $0<\gamma<1$. The smooth obstruction, positive-flat existence, and approximation ingredients are classical; the contribution is their explicit formulation and synthesis, with finite-order statements and complete reductions.
\end{abstract}

\noindent\textbf{2020 Mathematics Subject Classification.} 26E10, 52C35.

\noindent\textbf{Keywords.} finite jets; hyperplane arrangements; Denjoy--Carleman classes; quasianalyticity; sign distributions; polynomial sign representation.

\section{Introduction}

A Boolean truth table is finite. Requiring the same truth value throughout an entire orthant of $\R^n$ is not. This distinction becomes sharp when the Boolean value is extracted from the sign of a smooth real-valued function.

For example, the polynomial
\[
F(x,y)=xy
\]
is positive in the first and third quadrants and negative in the other two. With positive interpreted as $1$ and negative as $0$, it realizes equivalence (XNOR), an affine Boolean function over $\Ftwo$. By contrast, the AND sign pattern is positive only in the positive-positive quadrant and negative in the other three. The nonsmooth function $\min(x,y)$ has exactly that sign pattern, but the results below show that any \emph{smooth} realization of AND on all four quadrants must be flat at the origin. In a quasianalytic class this is impossible for a nonzero function; in a non-quasianalytic class it is possible by using flat one-sided selectors.

The first purpose of this paper is to isolate the local mechanism behind that phenomenon. It is not coordinate division, resolution of singularities, or a special feature of one ultradifferentiable class. It is the first nonzero homogeneous Taylor term. If a smooth function has a nonzero finite jet at the common corner, its leading polynomial already inherits the orthant signs in a weak sense. Polynomial factorization across coordinate hyperplanes then forces an affine parity law. The last factor-multiplicity step is elementary and has clear antecedents; the point of the finite-jet formulation is that it transports that obstruction from polynomials to arbitrary smooth germs and makes the subsequent regularity classification immediate.

The second purpose is to carry that mechanism to a finite affine hyperplane arrangement. If a smooth function is nonzero off the arrangement and has no flat point on it, the local leading polynomial at each point fixes a crossing parity for every incident hyperplane. Those parities propagate along each connected affine hyperplane and hence determine the sign on every chamber.

The third purpose is to combine these smooth statements with the classical Denjoy--Carleman theorem. We use throughout one compactwise Roumieu convention,
\[
 |D^\alpha f(x)|\le C_K\rho_K^{|\alpha|}|\alpha|!M_{|\alpha|},
 \qquad x\in K\Subset U,
\]
and no derivative-stability assumption. In this normalization the quasianalytic criterion is
\[
 \sum_{k=0}^{\infty}\frac{M_k}{(k+1)M_{k+1}}=\infty.
\]
On the divergent side, quasianalyticity rules out the flatness required by a non-affine orthant law. On the convergent side, a classical compactly supported cutoff yields a one-sided selector, from which arbitrary Boolean laws are assembled without cancellation.

\paragraph{Prior work and contribution boundary.}
The $C^\infty$ arrangement obstruction belongs to the signature-flatness theory of Broglia and Tognoli~\cite{BrogliaTognoli1989}; Remark~\ref{rem:BT} gives the precise reduction from the hypotheses used here, including the homological point in their theorem. Their factorization analysis also contains the parity-of-multiplicity mechanism underlying Lemma~\ref{thm:T1}, and their examples already use a first nonzero Taylor coefficient to detect incompatible crossing parities. Thus neither the smooth obstruction nor the leading-term idea is claimed as new.

The broader historical setting includes distributions of signs and signatures~\cite{Fortuna1991,AcquistapaceBroglia1992,AcquistapaceBroglia1993Correction,AcquistapaceBrogliaFortuna1994,AcquistapaceBroglia1994}, the Varchenko--Gelfand algebra of chamber functions and its later arrangement-theoretic extensions~\cite{VarchenkoGelfand1987,Moseley2017,BjornerEtAl1999}, and the classical Denjoy--Carleman theory of quasianalyticity, cutoffs, and positive-flat functions~\cite{Rainer2022,Langenbruch2003,Thilliez2008}. Stronger division and monomialization frameworks are also available~\cite{BierstoneMilman2004,BelottoBierstoneChow2018,BelottoBierstone2023}, but they are not needed for the elementary proofs below.

Exact Boolean sign operations also have a classical constructive theory in the $R$-function literature. Shapiro's survey~\cite[Sections~3.3 and 3.6]{Shapiro2007} describes root-based conjunction and disjunction and their finite-differentiability upgrades. These supply direct examples at the boundary of the finite-jet statements; they are credited explicitly in Section~\ref{sec:rfunctions}.

The note combines explicit finite-order rigidity, a self-contained proof for affine arrangements, and exact realization classifications with constructive formulas and chamber counts. Section~\ref{sec:quantitative} adds quantitative consequences: Taylor estimates turn vanishing jets into corner-margin bounds, while classical rational approximation, Sherstov's majority-conjunction lower bound~\cite{Sherstov2013}, and Bernstein--Jackson inequalities~\cite{KopotunPrymakShevchuk2018} give a fixed-dimensional degree classification after excluding the walls. These are derived consequences, with no claim of historical firstness or new approximation machinery. They distinguish exact realization near the common corner from sign correctness at a positive distance from the walls.

\paragraph{Organization.}
Section~\ref{sec:setup} fixes the sign and Denjoy--Carleman conventions. Section~\ref{sec:polynomial} proves the weak-sign polynomial lemma. Section~\ref{sec:finitejet} establishes finite-jet rigidity and the smooth arrangement theorem. Section~\ref{sec:dc} constructs one-sided selectors and proves the whole-orthant dichotomy. Section~\ref{sec:arrangements} gives the arrangement classification and exact counts. Section~\ref{sec:quantitative} derives corner-margin and polynomial-degree bounds. Section~\ref{sec:scope} records counterexamples and comparisons with $R$-functions and normal-crossing results.

\section{Definitions, conventions, and classical input}\label{sec:setup}

\subsection{Orthants and Boolean sign laws}

For $\alpha=(\alpha_1,\dots,\alpha_n)\in\bits^n$, put
\[
 \eps_i(\alpha)=2\alpha_i-1\in\{-1,+1\}
\]
and define the open orthant
\[
 \mathcal{O}_\alpha
 =\{x\in\R^n:\eps_i(\alpha)x_i>0\text{ for every }i\}.
\]
A function $F$ defined on an open set $U$ meeting all these orthants \emph{realizes} a Boolean law $\Theta:\bits^n\to\bits$ on the orthants in $U$ if $F$ is nonzero on their intersections with $U$ and
\[
 \sign F(x)=2\Theta(\alpha)-1,
 \qquad x\in\mathcal{O}_\alpha\cap U.
\]
Local realization at the common corner means realization on a centered ball contained in $U$; global realization means $U=\R^n$. We call $\Theta$ \emph{affine} if
\[
 \Theta(\alpha)=b\oplus e_1\alpha_1\oplus\cdots\oplus e_n\alpha_n,
 \qquad b,e_i\in\bits,
\]
where $\oplus$ denotes addition in $\Ftwo$. There are $2^{n+1}$ affine Boolean functions. When $n=1$, these are all four unary Boolean functions.

A smooth function is \emph{flat at $p$} if $D^\alpha F(p)=0$ for every multi-index $\alpha$.

\begin{definition}[Taylor-quasianalytic class]\label{def:taylorquasi}
For the two generic corollaries below, a smooth function class is called \emph{Taylor-quasianalytic} if, at every point of its domain, a function whose full Taylor jet vanishes has zero germ there. For a vector space of germs, this zero-jet implication is equivalent to injectivity of the Taylor map. Only the zero-jet implication is used in the two generic corollaries below. No division, composition, or sheaf axiom is intended by this shorthand. The divergent Denjoy--Carleman classes used later have this property by the Denjoy--Carleman theorem.
\end{definition}

\subsection{Affine arrangements}

Let
\[
 \cA=\{H_1,\dots,H_m\},\qquad H_j=\{x\in\R^n:\ell_j(x)=0\},
\]
be a finite family of distinct proper affine hyperplanes, where each nonconstant affine defining form $\ell_j$ is fixed once and for all. A \emph{chamber} is a connected component of
\[
 \R^n\setminus\bigcup_{j=1}^m H_j.
\]
We write $\Ch(\cA)$ for the chamber set. On a chamber $C$, every $\sign\ell_j$ is constant. A chamber labeling $h:\Ch(\cA)\to\{\pm1\}$ is a \emph{signed parity law} if
\begin{equation}\label{eq:signed-parity}
 h(C)=c\prod_{j=1}^m \sign(\ell_j|_C)^{e_j},
 \qquad c\in\{\pm1\},\quad e_j\in\bits.
\end{equation}
The word ``signed'' is deliberate: when $c=-1$, this is not literally a group character, and the realized chamber sign vectors need not form a group.

\subsection{One Denjoy--Carleman convention}

Throughout Sections~\ref{sec:dc}--\ref{sec:arrangements}, a \emph{normalized weight} means a sequence $M=(M_k)_{k\ge0}$ such that
\[
 M_0=1,
 \qquad M_k>0,
\]
with $M$ nondecreasing and log-convex. For an open set $U\subseteq\R^n$, define the compactwise Roumieu class $\mathcal{E}^{\{M\}}(U)$ by requiring that for every compact $K\Subset U$ there exist $C_K,\rho_K>0$ such that
\begin{equation}\label{eq:roumieu}
 |D^\alpha f(x)|
 \le C_K\rho_K^{|\alpha|}|\alpha|!M_{|\alpha|}
 \qquad (x\in K).
\end{equation}
This convention keeps the factorial outside $M$. To compare with the common total-weight convention, set
\[
 W_k=k!M_k.
\]
Then $W$ is a weight sequence in the sense used in \cite{Rainer2022}: it is positive and log-convex, $W_0=1\le W_1$, and $W_k^{1/k}\to\infty$. Its quotient sequence satisfies
\[
 \frac{W_{k+1}}{W_k}
 =(k+1)\frac{M_{k+1}}{M_k},
\]
so the Denjoy--Carleman series is precisely
\begin{equation}\label{eq:dcseries}
 \sum_{k=0}^{\infty}\frac{M_k}{(k+1)M_{k+1}}.
\end{equation}
The Denjoy--Carleman theorem says that $\mathcal{E}^{\{M\}}$ is quasianalytic exactly when \eqref{eq:dcseries} diverges; see \cite[Theorem~3.6]{Rainer2022}. In the convergent case, nontrivial compactly supported functions exist; the cutoff form we use is supplied, for example, by \cite[Corollary~3.12]{Rainer2022}.

We will need only elementary closure properties in addition to this classical input.

\begin{lemma}[Elementary closure]\label{lem:closure}
For a normalized weight $M$, the class $\mathcal{E}^{\{M\}}$ is closed under finite sums, finite products, translations, and precomposition with affine linear functionals. For functions on $\R$ it is also closed under the primitive $G(t)=\int_0^t g(s)\,ds$.
\end{lemma}

\begin{proof}
Only products and primitives need comment. Log-convexity with $M_0=1$ implies $M_jM_{k-j}\le M_k$. Combining this with Leibniz' rule and the bounds \eqref{eq:roumieu} gives the same form of estimate with a larger radius constant.

For a one-variable $g\in\mathcal{E}^{\{M\}}(\R)$, define $G(t)=\int_0^t g(s)\,ds$. On a compact interval, $G$ itself is bounded. For $k\ge1$, $G^{(k)}=g^{(k-1)}$, and monotonicity of $M$ gives
\[
 |G^{(k)}|
 \le C\rho^{k-1}(k-1)!M_{k-1}
 \le C'({\rho}')^k k!M_k
\]
for suitable enlarged constants. For an affine functional $\ell(x)=a\cdot x+b$ and a one-variable function $g$, the chain rule gives
\[
 D^\alpha(g\circ\ell)(x)=a^\alpha g^{(|\alpha|)}(\ell(x)).
\]
The image of a compact set is compact, and the factors from $a$ are absorbed into the radius constant. Translations are the one-variable special case of slope one.
\end{proof}

\section{A weak-sign polynomial lemma}\label{sec:polynomial}

The finite-jet argument rests on an elementary fact: even a \emph{weak} chamber sign prescription for a nonzero polynomial has parity form.

\begin{lemma}[Polynomial weak-sign lemma]\label{thm:T1}
Let $\cA=\{H_1,\dots,H_m\}$ be a finite family of distinct proper affine hyperplanes in $\R^n$, with fixed defining forms $\ell_j$. Let $P$ be a nonzero real polynomial. Suppose numbers $h_C\in\{\pm1\}$ are assigned to the chambers $C\in\Ch(\cA)$ and satisfy
\begin{equation}\label{eq:weaksign}
 h_C P(x)\ge0
 \qquad (x\in C).
\end{equation}
Then there are unique $c\in\{\pm1\}$ and $e_1,\dots,e_m\in\bits$ such that
\begin{equation}\label{eq:polyparity}
 h_C=c\prod_{j=1}^m\sign(\ell_j|_C)^{e_j}
 \qquad(C\in\Ch(\cA)).
\end{equation}
The polynomial is allowed to vanish at points inside a chamber.
\end{lemma}

\begin{proof}
For each $j$, let $a_j$ be the largest exponent for which $\ell_j^{a_j}$ divides $P$, and write
\[
 P=\ell_j^{a_j}Q_j,
 \qquad \ell_j\nmid Q_j.
\]
The restriction of $Q_j$ to $H_j$ is therefore a nonzero polynomial on $H_j$ and cannot vanish on a relatively open subset of $H_j$.

Consider two chambers sharing a facet contained in $H_j$. In the relative interior of that facet choose a point $p$ outside every other hyperplane and with $Q_j(p)\ne0$. In a sufficiently small ball about $p$, the sign of $Q_j$ is fixed, while $\ell_j$ changes sign when $H_j$ is crossed. Hence the sign of $P$ changes by the factor $(-1)^{a_j}$. Wherever $P\ne0$, the weak inequality \eqref{eq:weaksign} forces its sign to equal the prescribed chamber sign. Thus every adjacency across $H_j$ has the same ratio
\[
 \frac{h_{C'}}{h_C}=(-1)^{a_j}.
\]
Put $e_j\equiv a_j\pmod2$.

Now divide the chamber label by $\prod_j\sign(\ell_j)^{e_j}$. The quotient is unchanged across every chamber adjacency. The chamber adjacency graph is connected: choose interior points of two chambers and perturb a segment, or use a two-segment polygonal path, so that it meets the hyperplanes transversely and never meets two at the same time. The quotient is therefore a single constant $c\in\{\pm1\}$, proving \eqref{eq:polyparity}.

For uniqueness, fix $j$ and choose a point of $H_j$ outside all other hyperplanes. The two local chambers on its two sides differ only in the sign of $\ell_j$, so their label ratio determines $e_j$. One chamber then determines $c$.
\end{proof}

\begin{remark}
The content of Lemma~\ref{thm:T1} is elementary polynomial factorization plus chamber adjacency. We use it as a transparent engine for the smooth results; we do not attach an independent novelty claim to it.
\end{remark}

\section{Finite-jet rigidity}\label{sec:finitejet}

\subsection{Orthants}

\begin{theorem}[Finite-jet rigidity]\label{thm:T2}
Let $F$ be $C^\infty$ on a neighborhood of the origin in $\R^n$. Assume that in some ball $B$ centered at the origin,
\[
 F(x)\ne0
 \qquad\text{whenever every coordinate of $x$ is nonzero.}
\]
Because each $B\cap\mathcal O_\alpha$ is connected, continuity makes the sign of $F$ constant there; write the resulting local orthant sign law as $h:\{\pm1\}^n\to\{\pm1\}$. If at least one derivative of $F$ at the origin is nonzero, then there are unique $c\in\{\pm1\}$ and $e\in\bits^n$ such that
\begin{equation}\label{eq:orthparity}
 h(s)=c\prod_{i=1}^n s_i^{e_i}.
\end{equation}
Equivalently, every smooth realization of a non-affine Boolean orthant law is flat at the origin.
\end{theorem}

\begin{proof}
Let $d$ be the least order of a nonzero derivative of $F$ at the origin; $d=0$ is allowed. The first nonzero homogeneous Taylor term is
\[
 P_d(x)=\sum_{|\alpha|=d}\frac{D^\alpha F(0)}{\alpha!}x^\alpha,
 \qquad P_d\ne0.
\]
For any $v\in\R^n$ with no zero coordinate, Taylor's formula along the ray $tv$, $t\downarrow0$, gives
\[
 F(tv)=t^dP_d(v)+o(t^d).
\]
By the nonvanishing hypothesis and connectedness just noted, $F(tv)$ has sign $h(\sign v)$ for every sufficiently small $t>0$,
\[
 h(\sign v)\,P_d(v)\ge0.
\]
Thus the nonzero polynomial $P_d$ satisfies a weak sign prescription on the coordinate-hyperplane arrangement. Lemma~\ref{thm:T1} gives \eqref{eq:orthparity} and its uniqueness.

Under the identification $s_i=2\alpha_i-1$, a sign law of the form \eqref{eq:orthparity} is equivalent to an affine Boolean function $b\oplus e_1\alpha_1\oplus\cdots\oplus e_n\alpha_n$. The contrapositive gives the flatness statement.
\end{proof}

\begin{remark}[Finite regularity]
The proof uses only the first nonzero Taylor order. If the first nonzero jet has order $d$, then a $C^d$ Taylor expansion is enough. The theorem does not assert that nonflatness is necessary for an affine law: flat functions can also realize affine laws.
\end{remark}

\begin{theorem}[Finite-order orthant rigidity]\label{thm:finiteC}
Let $d\ge 0$ and let $F$ be $C^d$ near the origin, nonzero whenever all coordinates are nonzero in some centered ball. If every derivative of order $<d$ vanishes at the origin and at least one derivative of order $d$ is nonzero, then the local orthant sign law is affine parity.
\end{theorem}

\begin{proof}
The order-$d$ Taylor expansion supplies the same nonzero leading homogeneous polynomial used in Theorem~\ref{thm:T2}; Lemma~\ref{thm:T1} then applies verbatim.
\end{proof}

\begin{corollary}[Quasianalytic orthant rigidity]\label{cor:quasiorth}
Let $F$ belong to a Taylor-quasianalytic class on a neighborhood of the origin and suppose it realizes a whole local orthant law there. Then that Boolean law is affine.
\end{corollary}

\begin{proof}
If every derivative of $F$ vanished at the origin, the defining zero-jet property would make the germ of $F$ at the origin zero, contradicting nonvanishing on every local open orthant. Apply Theorem~\ref{thm:T2}.
\end{proof}

\subsection{Affine hyperplane arrangements}

\begin{theorem}[No-flat-point arrangement rigidity]\label{thm:T3}
Let $\cA=\{H_1,\dots,H_m\}$ be a finite family of distinct proper affine hyperplanes in $\R^n$. Let $F\in C^\infty(\R^n)$ be nonzero on $\R^n\setminus\bigcup_jH_j$. Assume that $F$ is not flat at any point of $\bigcup_jH_j$. Then there are unique $c\in\{\pm1\}$ and $e_1,\dots,e_m\in\bits$ such that
\begin{equation}\label{eq:arrparity}
 \sign F(x)
 =c\prod_{j=1}^m\sign\ell_j(x)^{e_j},
 \qquad x\notin\bigcup_jH_j.
\end{equation}
\end{theorem}

\begin{proof}
Fix $p\in\bigcup_jH_j$. Choose a ball $B_p$ small enough that every hyperplane meeting $B_p$ contains $p$. After translating $p$ to the origin, the incident hyperplanes form a central arrangement. Let $P_p$ be the first nonzero homogeneous Taylor polynomial of $F$ at $p$.

For every direction $v$ outside the incident hyperplanes,
\[
 F(p+tv)=t^{d(p)}P_p(v)+o(t^{d(p)}),
 \qquad t\downarrow0.
\]
The function $F$ has a fixed sign on each local chamber because it is continuous and nonzero there. Hence $P_p$ inherits those signs weakly. Lemma~\ref{thm:T1} gives, for every hyperplane $H_j$ through $p$, a local crossing parity $e_j(p)\in\bits$.

We claim that $e_j(p)$ is locally constant as $p$ varies in $H_j$. Fix $p\in H_j$ and its ball $B_p$. For $p'\in H_j$ sufficiently close to $p$, choose the local ball $B_{p'}$ contained in $B_p$ and avoiding all hyperplanes not containing $p'$. The relatively open set $H_j\cap B_{p'}$ contains a point $q$ outside every other hyperplane: the excluded intersections form a finite union of proper affine subspaces of $H_j$. In dimension one, $H_j$ is a point and distinct hyperplanes do not meet it. A sufficiently small transverse pair near $q$, on opposite sides of $H_j$, lies inside both balls and meets no other hyperplane. The ratio of the two actual signs of $F$ is both $(-1)^{e_j(p)}$ and $(-1)^{e_j(p')}$. Thus the bits agree. Since $H_j$ is connected, $e_j(p)$ is a single global bit $e_j$.

Now consider the global chamber adjacency graph. Across a generic facet contained in $H_j$, the sign of $F$ changes by $(-1)^{e_j}$. The same is true of the product $\prod_j\sign(\ell_j)^{e_j}$. Their quotient is therefore unchanged across every chamber adjacency. As in Lemma~\ref{thm:T1}, the chamber adjacency graph is connected, so the quotient is a global constant $c\in\{\pm1\}$. This proves \eqref{eq:arrparity}. Uniqueness follows by crossing a generic point of each $H_j$.
\end{proof}

\begin{remark}[Earlier signature-flatness theorem]\label{rem:BT}
Theorem~\ref{thm:T3} is an elementary affine-arrangement case of the signature-flatness theory of Broglia and Tognoli~\cite[Theorem~4.2]{BrogliaTognoli1989}. To compare hypotheses, set
\[
 G=F\prod_{j=1}^m\ell_j^2.
\]
Then $G$ has zero set exactly $\bigcup_jH_j$ and the same chamber signs as $F$. If $F$ has a first nonzero Taylor term at a point of the arrangement, the first nonzero term of $G$ is its product with a nonzero polynomial, so nonflatness is preserved. For $n\ge2$ their homological condition is automatic on $\R^n$ because $H_{n-1}(\R^n;\mathbb Z_2)=0$; the hyperplane components admit global affine equations and satisfy the required connectivity. The case $n=1$ is immediate from the one-dimensional crossing picture. Their admissible signatures reduce, for hyperplanes, to the signed parity laws in \eqref{eq:signed-parity}. This remark records the historical implication explicitly; the proof above is retained because it isolates the finite-jet mechanism needed for Theorem~\ref{thm:finitearr}.
\end{remark}

\begin{theorem}[Uniform finite-order arrangement rigidity]\label{thm:finitearr}
Let $N\ge0$, let $\cA=\{H_1,\dots,H_m\}$ be a finite family of distinct proper affine hyperplanes in $\R^n$, and let $F\in C^N(\R^n)$ be nonzero on $\R^n\setminus\bigcup_jH_j$. Assume that for every $p\in\bigcup_jH_j$ the $N$-jet of $F$ at $p$ is nonzero; equivalently, for each such $p$ there is a multi-index $\alpha$ with $|\alpha|\le N$ and $D^\alpha F(p)\ne0$. Then the chamber sign law of $F$ is a unique signed parity law of the form \eqref{eq:arrparity}.
\end{theorem}

\begin{proof}
At each $p$ on the arrangement choose the least $d(p)\le N$ for which some derivative of order $d(p)$ is nonzero. The $C^{d(p)}$ Taylor expansion gives a nonzero leading homogeneous polynomial and therefore the same local crossing parities as in the proof of Theorem~\ref{thm:T3}. The propagation of those parities along a hyperplane, and then across the connected chamber adjacency graph, uses only continuity and nonvanishing off the arrangement. Thus the remainder of that proof applies unchanged.
\end{proof}

\begin{remark}[A finite-order obstruction]
Equivalently, if a $C^N$ function is nonzero off the arrangement and realizes a chamber labeling that is not signed parity, then at some point of the arrangement its entire $N$-jet must vanish. The smooth theorem is a stronger all-orders analogue: it yields a single flat point. The finite-order statements alone do not imply that conclusion by letting $N$ tend to infinity, because the point at which the $N$-jet vanishes may depend on $N$.
\end{remark}

\begin{corollary}[Quasianalytic arrangement rigidity]\label{cor:quasiarr}
Let $F$ be a nonzero globally defined member of a Taylor-quasianalytic class on $\R^n$, and suppose $F$ is nonzero off a finite affine arrangement. Then its chamber sign law has the form \eqref{eq:arrparity}.
\end{corollary}

\begin{proof}
At a flat point, the defining zero-jet property would make the germ of $F$ there zero. Every neighborhood of a point on a finite hyperplane arrangement contains points off the arrangement, contradicting the assumed nonvanishing there. Thus $F$ has no flat point on the arrangement, and Theorem~\ref{thm:T3} applies.
\end{proof}

\section{Denjoy--Carleman realization on orthants}\label{sec:dc}

The finite-jet theorem supplies the rigidity half of the classification. The converse direction in a non-quasianalytic class is constructive.

\begin{lemma}[One-sided selector]\label{thm:T4}
Let $M$ be a normalized weight. If
\begin{equation}\label{eq:nonq}
 \sum_{k=0}^{\infty}\frac{M_k}{(k+1)M_{k+1}}<\infty,
\end{equation}
then there exists $\psi\in\mathcal{E}^{\{M\}}(\R)$ such that
\[
 \psi(t)=0\quad(t\le0),
 \qquad
 \psi(t)>0\quad(t>0).
\]
\end{lemma}

\begin{proof}
With $W_k=k!M_k$, condition \eqref{eq:nonq} is the non-quasianalytic Denjoy--Carleman condition $\sum W_k/W_{k+1}<\infty$. By the classical cutoff theorem, choose a nonzero real-valued compactly supported function $f$ in the corresponding Roumieu class; see \cite[Theorem~3.6 and Corollary~3.12]{Rainer2022}. Translate $f$ so that the left endpoint of its support is $0$. Then
\[
 f(t)=0\quad(t\le0),
\]
and every interval $(0,t)$, $t>0$, contains a point at which $f$ is nonzero: otherwise $f$ would vanish on a right-neighborhood of $0$, contradicting that $0$ is the left endpoint of $\operatorname{supp}f$.

By Lemma~\ref{lem:closure}, $f^2$ and its primitive belong to the same class. Define
\[
 \psi(t)=\int_0^t f(s)^2\,ds.
\]
For $t\le0$ the integral is zero. For $t>0$, the interval $(0,t)$ contains a point where $f\ne0$; continuity then gives an interval on which $f^2>0$, and therefore $\psi(t)>0$.
\end{proof}

\paragraph{Classical constructive antecedent.}
A stronger existence statement is classical. Rainer's Remark~3.13, citing work of Langenbruch, records positive-flat ultradifferentiable functions for prescribed closed sets in the non-quasianalytic setting~\cite[Remark~3.13]{Rainer2022}; see also~\cite{Langenbruch2003}. Taking the closed set to be the arrangement and changing the sign of such a function independently on each chamber already gives arbitrary chamber labelings, since flatness makes the chamberwise pieces glue with zero jet on every wall. The selector construction above is retained because it gives an explicit formula in the defining affine forms and makes the later counting argument transparent.

\begin{definition}
For a normalized weight $M$, let $\mathcal{R}_M(n)$ be the set of Boolean functions $\Theta:\bits^n\to\bits$ realized by some globally defined $F\in\mathcal{E}^{\{M\}}(\R^n)$ that is nonzero whenever all coordinates are nonzero.
\end{definition}

\begin{theorem}[Whole-orthant classification]\label{thm:T5}
Let $M$ be a normalized weight. Then
\[
 \mathcal{R}_M(n)=
 \begin{cases}
 \{\text{affine Boolean functions}\},
 &\displaystyle\sum_{k=0}^{\infty}\frac{M_k}{(k+1)M_{k+1}}=\infty,\\[1.1em]
 \{\text{all Boolean functions on }\bits^n\},
 &\displaystyle\sum_{k=0}^{\infty}\frac{M_k}{(k+1)M_{k+1}}<\infty.
 \end{cases}
\]
\end{theorem}

\begin{proof}
Assume first that the series diverges. By the Denjoy--Carleman theorem, $\mathcal{E}^{\{M\}}$ is quasianalytic. If $F$ realizes a whole-orthant law, then $F$ is not the zero function in any neighborhood of the origin, so its Taylor jet at the origin cannot vanish identically. Theorem~\ref{thm:T2} implies that the law is affine. Conversely, every affine law is realized by the polynomial
\[
 F(x)=c\prod_{i=1}^n x_i^{e_i}
\]
for suitable $c\in\{\pm1\}$ and $e_i\in\bits$, and every polynomial lies in the class.

Now assume the series converges and choose the selector $\psi$ from Lemma~\ref{thm:T4}. For $\alpha\in\bits^n$, define
\[
 \Psi_\alpha(x)
 =\prod_{i=1}^n\psi\bigl(\eps_i(\alpha)x_i\bigr).
\]
For an arbitrary Boolean function $\Theta$, set
\begin{equation}\label{eq:booleanselector}
 F_\Theta(x)
 =\sum_{\alpha\in\bits^n}
 \bigl(2\Theta(\alpha)-1\bigr)\Psi_\alpha(x).
\end{equation}
By Lemma~\ref{lem:closure}, $F_\Theta\in\mathcal{E}^{\{M\}}(\R^n)$. If $x\in\mathcal{O}_\alpha$, then $\Psi_\alpha(x)>0$, while for every $\beta\ne\alpha$ at least one factor of $\Psi_\beta(x)$ is zero. Hence
\[
 \sign F_\Theta(x)=2\Theta(\alpha)-1.
\]
Thus every Boolean function is realized.
\end{proof}

\begin{corollary}[Gevrey dichotomy]\label{cor:gevrey}
For the Gevrey class $G^s$, $s\ge1$, with derivative growth $C\rho^k(k!)^s$, the whole-orthant laws are exactly the affine Boolean functions when $s=1$, and all Boolean functions when $s>1$.
\end{corollary}

\begin{proof}
In the normalization \eqref{eq:roumieu}, take $M_k=(k!)^{s-1}$. Then
\[
 \frac{M_k}{(k+1)M_{k+1}}=(k+1)^{-s}.
\]
The corresponding series diverges exactly at $s=1$ for $s\ge1$. Apply Theorem~\ref{thm:T5}.
\end{proof}

\begin{remark}
For $n=1$, the two sides of Theorem~\ref{thm:T5} contain the same four Boolean functions because every unary Boolean function is affine. The strict expressive difference begins at arity $n\ge2$.
\end{remark}

\section{Affine arrangements: classification and exact counts}\label{sec:arrangements}

The orthant construction extends directly once orthants are replaced by the realized sign cells of a finite affine arrangement.

\begin{theorem}[Arrangement classification]\label{thm:T6}
Let $M$ be a normalized weight. Let $\cA$ be an arrangement of $m$ distinct proper affine hyperplanes in $\R^n$, and write $r=|\Ch(\cA)|$.

If the series \eqref{eq:dcseries} diverges, the chamber sign laws realized by globally defined $F\in\mathcal{E}^{\{M\}}(\R^n)$ that are nonzero off the arrangement are exactly the signed parity laws \eqref{eq:signed-parity}; there are exactly $2^{m+1}$ such laws.

If the series \eqref{eq:dcseries} converges, every one of the $2^r$ chamber labelings is realizable by such a function.
\end{theorem}

\begin{proof}
On the quasianalytic side, Corollary~\ref{cor:quasiarr} gives the signed-parity form. Every choice of $(c,e_1,\dots,e_m)$ is realized by the polynomial
\[
 c\prod_{j=1}^m\ell_j(x)^{e_j}.
\]
These $2^{m+1}$ restrictions are distinct. Indeed, for each $H_j$ choose a generic point of $H_j$ outside all other hyperplanes. The adjacent chamber pair at that point differs only in the sign of $\ell_j$, so its label ratio recovers $e_j$; one chamber then recovers $c$.

On the non-quasianalytic side, let $\psi$ be the selector from Lemma~\ref{thm:T4}. For each chamber $C$, let
\[
 s_j(C)=\sign(\ell_j|_C)
\]
and define
\[
 \Psi_C(x)=\prod_{j=1}^m\psi\bigl(s_j(C)\ell_j(x)\bigr).
\]
Distinct chambers have distinct sign vectors: a nonempty fixed sign cell is an intersection of open half-spaces and is therefore convex and connected. Hence if $x$ lies in a chamber $D$, then $\Psi_D(x)>0$ and $\Psi_C(x)=0$ for every $C\ne D$.

Given any labeling $h:\Ch(\cA)\to\{\pm1\}$, set
\[
 F_h(x)=\sum_{C\in\Ch(\cA)} h(C)\Psi_C(x).
\]
Elementary closure gives $F_h\in\mathcal{E}^{\{M\}}(\R^n)$, and on chamber $D$ the sum reduces to $h(D)\Psi_D(x)$. Thus $F_h$ has the prescribed sign and is nonzero off the arrangement.
\end{proof}

\begin{example}[Three concurrent lines]
Take three distinct lines through the origin in $\R^2$, for example
\[
 x=0,\qquad y=0,\qquad x-y=0.
\]
There are $r=6$ chambers. A full quasianalytic normalized Roumieu class $\mathcal{E}^{\{M\}}$ therefore realizes $2^{3+1}=16$ signed-parity chamber laws, whereas a full non-quasianalytic such class realizes all $2^6=64$ labelings.
\end{example}

\begin{remark}[When the counts coincide]
Each added distinct hyperplane splits at least one existing chamber, so $r\ge m+1$. The numerical gap is strict exactly when $r>m+1$. For $m$ parallel hyperplanes one has $r=m+1$, so the two counts coincide. Thus the quasianalytic/non-quasianalytic distinction is structural, but its cardinality shadow depends on the arrangement.
\end{remark}

\begin{table}[t]
\centering
\caption{The classification in the common notation of this paper. All exact realizability counts refer to the full normalized Roumieu classes $\mathcal{E}^{\{M\}}$.}
\label{tab:summary}
\begin{tabularx}{\textwidth}{@{}>{\raggedright\arraybackslash}p{0.24\textwidth}X>{\raggedright\arraybackslash}p{0.24\textwidth}@{}}
\toprule
Setting & Mechanism & Sign laws \\
\midrule
Smooth germ with a nonzero jet at the common orthant corner & Leading homogeneous Taylor polynomial and Lemma~\ref{thm:T1} & Affine parity only \\
$C^N$ function on an arrangement with nonzero $N$-jet at every wall point & Local leading terms plus parity propagation & Signed parity only \\
Quasianalytic $\mathcal{E}^{\{M\}}$ on orthants & Taylor injectivity excludes the flatness required by Theorem~\ref{thm:T2} & Exactly $2^{n+1}$ affine Boolean laws \\
Non-quasianalytic $\mathcal{E}^{\{M\}}$ on orthants & One-sided flat selector & All $2^{2^n}$ Boolean laws \\
Quasianalytic $\mathcal{E}^{\{M\}}$ on $m$ distinct hyperplanes & No flat points plus Theorem~\ref{thm:T3} & Exactly $2^{m+1}$ signed-parity laws \\
Non-quasianalytic $\mathcal{E}^{\{M\}}$ on an arrangement with $r$ chambers & Chamber selectors & All $2^r$ chamber labelings \\
\bottomrule
\end{tabularx}
\end{table}

\section{Margins and exclusion of the coordinate walls}\label{sec:quantitative}

The qualitative classification imposes correctness arbitrarily close to every wall. Two quantitative consequences clarify this requirement. Throughout this section $h:\{-1,+1\}^n\to\{-1,+1\}$, with positive meaning true, and $n$ is fixed. A non-affine law means one not of the form $c\prod_{i\in S}s_i$.

\subsection{Corner margins under derivative bounds}

\begin{proposition}[A normalized corner estimate]\label{prop:corner}
Let $k\ge0$, $0<\alpha\le1$, and let $F\in C^{k,\alpha}(B_R(0))$ realize a non-affine orthant law. Suppose that $D^kF$ has H\"older seminorm at most $H$, using the Euclidean multilinear operator norm (absolute value for $k=0$). Then its entire $k$-jet at zero vanishes and
\begin{equation}\label{eq:cornerbound}
 |F(x)|\le H\frac{\Gamma(\alpha+1)}{\Gamma(k+\alpha+1)}\|x\|^{k+\alpha}
 \qquad(x\in B_R(0)).
\end{equation}
In particular, if $\sqrt n\delta<R$ and $h(s)F(\delta s)\ge\gamma$ for all sign vectors $s$, then
\[
 \gamma\le H\frac{\Gamma(\alpha+1)}{\Gamma(k+\alpha+1)}(\sqrt n\delta)^{k+\alpha}.
\]
\end{proposition}

\begin{proof}
Any least nonzero derivative of order at most $k$ would contradict Theorem~\ref{thm:finiteC}. Thus all these derivatives vanish. For $k=0$, \eqref{eq:cornerbound} is the H\"older estimate from $F(0)=0$. For $k\ge1$, write $x=rv$, $\|v\|=1$, and use integral Taylor remainder:
\[
 |F(rv)|\le\frac{H}{(k-1)!}\int_0^r(r-t)^{k-1}t^\alpha\,dt
 =H\frac{\Gamma(\alpha+1)}{\Gamma(k+\alpha+1)}r^{k+\alpha}.
\]
Substitute $r=\sqrt n\delta$ for the margin conclusion.
\end{proof}

\begin{remark}[Smooth derivative budgets]\label{rem:associated}
For a smooth exact non-affine realization, Theorem~\ref{thm:T2} gives flatness. If, on $B_R(0)$, fixed constants $A,B>0$ satisfy
\[
 \|D^jF(x)\|_{\mathrm{op}}\le A B^j j!M_j\qquad(j\ge0),
\]
Taylor remainder at every order gives
\begin{equation}\label{eq:associatedbound}
 |F(x)|\le A h_M(B\|x\|),\qquad h_M(t)=\inf_{j\ge0}M_jt^j.
\end{equation}
For $j=0$ use the amplitude bound; for $j\ge1$ the factorial cancels in the remainder. This is the usual associated-function estimate in the factorial-out convention; compare~\cite[Section~7]{Rainer2022}. Coordinate partial-derivative bounds of the form \eqref{eq:roumieu}, with fixed constants on the same ball (or on a fixed compact subball), imply the corresponding operator bounds after replacing $B$ by $B\sqrt n$. In the Gevrey case $M_j=(j!)^{s-1}$, $s>1$, optimization in $j$ gives $h_M(t)\le C\exp(-c t^{-1/(s-1)})$ for small $t>0$. A matching optimal-cutoff assertion for general weights requires additional assumptions and is not used here. These estimates concern spatial derivatives and do not assert failure of an optimization or learning procedure.
\end{remark}

\subsection{Polynomial signs at positive distance from the walls}

Put
\[
 Q_n=[-1,1]^n,\qquad
 D_{\delta,n}=\{x\in Q_n:|x_i|\ge\delta\text{ for all }i\},\quad 0<\delta<1.
\]
Let $d_h(\delta)$ be the least total degree of a real polynomial $P$ satisfying
\[
 h(\sign x)P(x)>0\qquad(x\in D_{\delta,n}).
\]
There is no amplitude or coefficient bound in this definition. For fixed $0<\gamma<1$, let $d^{\mathrm{bd}}_{h,\gamma}(\delta)$ be the least total degree subject instead to
\begin{equation}\label{eq:boundedmargin}
 \|P\|_{Q_n,\infty}\le1,\qquad
 h(\sign x)P(x)\ge\gamma\quad(x\in D_{\delta,n}).
\end{equation}
The amplitude bound covers the entire cube, including its excluded strips.

\begin{theorem}[Fixed-dimensional degree comparison]\label{thm:guardeddegree}
For fixed $n\ge1$ the following statements hold.
\begin{enumerate}[label=(\roman*)]
\item If $h(s)=c\prod_{i\in S}s_i$, then $d_h(\delta)=|S|$ for every $0<\delta<1$.
\item If $h$ is non-affine, then $d_h(\delta)=\Theta_h(\log(1/\delta))$ as $\delta\downarrow0$.
\item If $h$ is nonconstant and $0<\gamma<1$ is fixed, then
\[
 d^{\mathrm{bd}}_{h,\gamma}(\delta)=\Theta_{h,\gamma}(1/\delta)
 \qquad(\delta\downarrow0).
\]
In fact, $d^{\mathrm{bd}}_{h,\gamma}(\delta)\ge\gamma\sqrt{1-\delta^2}/\delta$.
\end{enumerate}
Constants may depend on the fixed dimension, law, and margin. No uniform growing-arity, coefficient-size, or evaluation-cost assertion is intended.
\end{theorem}

\begin{proof}
\emph{Parity.} The polynomial $c\prod_{i\in S}x_i$ proves the upper bound. On Boolean vertices, fix coordinates outside $S$. Every polynomial of degree less than $|S|$ has zero correlation with parity, by expanding into monomials and averaging over the remaining independent signs. Strict sign agreement would give positive correlation. This proves (i), including constants with $S=\varnothing$.

\emph{A logarithmic upper bound for every law.} Let $m=\lceil\log_2(1/\delta)\rceil$ and choose an integer $L\ge1$. Define
\[
 p(t)=\prod_{j=0}^m(t+2^{-j})^{2L},\quad
 a(t)=p(t)-p(-t),\quad q(t)=p(t)+p(-t),\quad r(t)=\frac{a(t)}{q(t)}.
\]
The two nonnegative polynomials $p(t)$ and $p(-t)$ have disjoint real zero sets, so $q>0$ and $|r|\le1$ on $\R$. For $2^{-m}\le t\le1$, some dyadic node $u=2^{-j}$ satisfies $u\le t\le2u$. Therefore
\[
 0\le\frac{p(-t)}{p(t)}
 =\prod_{j=0}^m\left|\frac{2^{-j}-t}{2^{-j}+t}\right|^{2L}
 \le3^{-2L},\qquad |r(t)-1|\le2\,3^{-2L}.
\]
Oddness gives the negative-interval estimate. This is a dyadic rational construction of Newman type; see also~\cite[Theorem~2.4]{Sherstov2013}.

For $u\in[-1,1]^n$, define the multilinear expectation extension
\[
 H_h(u)=\sum_{s\in\{-1,+1\}^n}h(s)\prod_{i=1}^n\frac{1+s_i u_i}{2}.
\]
It is bounded by one. If $|u_i-s_i|\le\eps$ for every coordinate, regard $H_h(u)$ as the expectation of $h$ under independent signs with means $u_i$. Each sign differs from $s_i$ with probability at most $\eps/2$, so the union bound gives
\begin{equation}\label{eq:extensionmargin}
 h(s)H_h(u)\ge1-n\eps.
\end{equation}
Choose $L$ with $2n3^{-2L}<1$. Substitution of $r(x_i)$ and multiplication by $\prod_iq(x_i)>0$ give a polynomial of total degree at most $2Ln(m+1)$, strictly correct on $D_{\delta,n}$. Hence $d_h(\delta)=O_n(\log(1/\delta))$.

\emph{The non-affine lower bound.} Every non-affine Boolean law has a non-affine two-dimensional face: otherwise every second Boolean difference of its $\Ftwo$-valued truth function would vanish identically, forcing its multilinear algebraic normal form to have degree at most one. Fix the other coordinates at the corresponding $\pm1$ values. Each of the eight non-affine binary laws is AND after input sign flips and possibly output negation, all of which preserve degree and the guarded square.

Choose an odd integer $N\le1/\delta$ with $N\asymp1/\delta$. In a polynomial for binary AND substitute
\[
 x=\frac1N\sum_{i=1}^N z_i,\qquad
 y=\frac1N\sum_{i=1}^N w_i,\qquad z_i,w_i\in\{-1,+1\}.
\]
Each magnitude is at least $1/N\ge\delta$. Thus the substituted polynomial sign-represents the conjunction of two independent $N$-variable majorities on a finite Boolean cube, without increasing degree. Sherstov's lower bound~\cite[Theorem~1.10]{Sherstov2013} is $\Omega(\log N)$. His convention uses $-1$ for true; input and output sign changes align the conventions. This proves (ii).

\emph{Bounded margin: lower bound.} A nonconstant law changes along some Boolean edge. Fix the other coordinates at that edge's $\pm1$ values. The resulting univariate slice $p$ is bounded by one on $[-1,1]$ and takes opposite values of magnitude at least $\gamma$ at $\pm\delta$. The mean value theorem gives $\xi\in(-\delta,\delta)$ with $|p'(\xi)|\ge\gamma/\delta$. Bernstein's inequality~\cite[Theorem~24.1]{KopotunPrymakShevchuk2018} gives, with $d=\deg P$,
\[
 \frac\gamma\delta\le |p'(\xi)|
 \le\frac d{\sqrt{1-\xi^2}}
 \le\frac d{\sqrt{1-\delta^2}}.
\]

\emph{Bounded margin: upper bound.} The clipped ramp
$v_\delta(t)=\max(-1,\min(1,t/\delta))$ is $1/\delta$-Lipschitz. Jackson's inequality yields a polynomial $u$ of degree at most $k$ with
\[
 \|u-v_\delta\|_{[-1,1],\infty}\le\frac C{k\delta}.
\]
For completeness, the trigonometric form~\cite[Theorem~21.1]{KopotunPrymakShevchuk2018} applies to the even periodic function $v_\delta(\cos\theta)$. Symmetrizing its approximant preserves the error and leaves only cosine terms; $\cos(j\theta)=T_j(\cos\theta)$ converts it to an algebraic polynomial of the same degree.
Choose a fixed error $\eta>0$ and $k\ge C/(\eta\delta)$. The polynomial $r=u/(1+\eta)$ is bounded by one on the whole interval and differs from $\sign t$ by at most $\eps=2\eta/(1+\eta)$ when $|t|\ge\delta$. Choose $\eta$ small enough that $n\eps\le1-\gamma$. Then $H_h(r(x_1),\ldots,r(x_n))$ satisfies \eqref{eq:boundedmargin}, by \eqref{eq:extensionmargin}, and has total degree at most $nk=O_{n,\gamma}(1/\delta)$. This proves (iii).
\end{proof}

\begin{remark}[Interpretation and endpoints]
Positive denominator clearing preserves signs but need not preserve a fixed margin after normalizing the polynomial on $Q_n$. This explains why (ii) and (iii) concern different resources. Constants have degree zero in both models; when $\gamma=1$, no nonconstant law satisfies \eqref{eq:boundedmargin}, since the polynomial would equal a constant on an open box and hence everywhere. The theorem is a synthesis of established approximation inputs. It does not follow from qualitative finite-jet rigidity alone, nor does it apply that rigidity to a polynomial required to be correct only on $D_{\delta,n}$.
\end{remark}

\section{Scope, counterexamples, and relation to nearby theory}\label{sec:scope}

\subsection{Real zero support is weaker than a normal-crossing divisor}

If one already has a factorization
\[
 F=u\prod_i x_i^{a_i}
\]
with a nonvanishing unit $u$, then the sign law is immediately the sign of $u$ times a parity product of coordinate signs. Normal-crossing and monomialization theories therefore provide important neighboring context; see \cite{BierstoneMilman2004,BelottoBierstone2023}. They should not, however, be substituted for the weaker hypothesis used in Theorems~\ref{thm:T3} and \ref{thm:T6}.

Indeed,
\begin{equation}\label{eq:weak-support-example}
 F(x,y)=xy(x^2+y^2)
\end{equation}
is real analytic and has real zero set exactly equal to the union of the two coordinate axes. Its sign off those axes is $\sign(xy)$. But after extracting the generic first-order factors $x$ and $y$, the residual factor $x^2+y^2$ still vanishes at the origin and is not a unit. Thus set-theoretic real zero support on an arrangement does not imply a unit-times-monomial factorization in the original coordinates.

This is why the arrangement theorem above is stated and proved directly at the sign level. Resolution and monomialization results involve stronger divisor hypotheses and/or transformations of source and target; they remain relevant background, not a replacement for the present argument.

\subsection{A common corner is essential for local orthant rigidity}

The connectedness of a domain alone does not replace the common-corner hypothesis in Theorem~\ref{thm:T2}. On the punctured plane $U=\R^2\setminus\{0\}$, define
\begin{equation}\label{eq:punctured-or}
 F(x,y)=x+y+\sqrt{x^2+y^2}.
\end{equation}
This function is real analytic on $U$. It is positive in the positive-positive quadrant and in both mixed-sign quadrants, while it is negative in the negative-negative quadrant. In a mixed-sign quadrant,
\[
 \sqrt{x^2+y^2}>|x+y|
\]
because $xy<0$, whereas in the negative-negative quadrant
\[
 \sqrt{x^2+y^2}<|x+y|=-x-y.
\]
Thus \eqref{eq:punctured-or} realizes OR on the four quadrants of the punctured plane. There is no contradiction with Theorem~\ref{thm:T2}: the origin, at which the finite-jet obstruction would live, has been removed.

\subsection{Classical R-functions and finite differentiability}\label{sec:rfunctions}

The preceding root formula is the classical $R$-disjunction. Changing the plus sign to a minus sign gives an $R$-conjunction; see~\cite[Section~3.3, equations (3.4)--(3.5)]{Shapiro2007}. In particular, for any integer $k\ge0$, put $r=(x^2+y^2)^{1/2}$ and
\[
 F_k(x,y)=r^k(x+y-r),\qquad F_k(0,0)=0.
\]
This realizes AND on the open quadrants. Its smooth angular factor and positive homogeneity of degree $k+1$ give $D^jF_k=O(r^{k+1-j})$ for $j\le k$; these derivatives extend by zero at the origin. The $k$th derivative is Lipschitz: for pairs at distance comparable to their radius use its $O(r)$ size bound, and for closer pairs integrate its bounded derivative along a segment away from zero. Thus $F_k\in C^{k,1}$ and its $k$-jet vanishes. It is not $C^{k+1}$: that regularity and homogeneity would make it its homogeneous Taylor polynomial of degree $k+1$, contrary to the non-affine sign law. On $(t,t)$, $t>0$, its value is $2^{k/2}(2-\sqrt2)t^{k+1}$. Hence the exponent $k+1$ in Proposition~\ref{prop:corner} for $\alpha=1$ cannot be increased in general. The construction is classical; the estimate states its relation to the finite-order obstruction.

\subsection{Analytic models and restricted domains}

A finite composition of real-analytic maps is real analytic on its domain. If a scalar output realizes an orthant law on a full neighborhood of the common corner, Corollary~\ref{cor:quasiorth} therefore forces affine parity. The hypothesis concerns independently variable coordinates throughout that neighborhood; it does not cover merely fitting finitely many Boolean representatives or a restricted set of reachable inputs.

Even an analytic formula can realize AND on every $D_{\delta,2}$ when it depends on $\delta$. For example,
\[
 F_\delta(x,y)=x+y-\sqrt{x^2+y^2+\delta^2}
\]
is globally real analytic. In the positive-positive guarded quadrant it is increasing in each coordinate, so its minimum is $(2-\sqrt3)\delta>0$. If $x\le-\delta$, then $F_\delta(x,y)<x+y-|y|\le-\delta$, and the other negative-input case is symmetric. At the common corner $F_\delta(0,0)=-\delta$, so it fails the positive-positive law sufficiently near zero. This example imposes no global amplitude normalization and does not contradict Theorem~\ref{thm:guardeddegree}. A nonsmooth alternative is $\min(x,y)=x-\max(0,x-y)$, which realizes AND exactly on every open quadrant.

\subsection{Smoothness alone does not imply rigidity}

The selector construction in \eqref{eq:booleanselector} gives $C^\infty$ realizations of every Boolean law in the non-quasianalytic case. For a non-affine law, Theorem~\ref{thm:T2} says that such a realization must be flat at the origin; the construction is consistent with exactly that requirement. The mathematical boundary is therefore not ``smooth versus nonsmooth,'' but rather the presence or absence of a nonzero finite jet at the relevant common boundary.

\subsection{Signatures and chamber functions}
The factor-parity mechanism in Lemma~\ref{thm:T1} is classical antecedent material. Broglia--Tognoli's Proposition~2.4 records how parity of maximal factor multiplicities controls the induced signature~\cite[Proposition~2.4]{BrogliaTognoli1989}, while the subsequent signature and sign-distribution literature studies broader realization and approximation questions~\cite{Fortuna1991,AcquistapaceBroglia1992,AcquistapaceBroglia1993Correction,AcquistapaceBrogliaFortuna1994,AcquistapaceBroglia1994}.

A distinct combinatorial tradition studies all functions that are constant on arrangement chambers. Varchenko and Gelfand introduced the algebra generated by hyperplane Heaviside functions~\cite{VarchenkoGelfand1987}; see also~\cite{Moseley2017,BjornerEtAl1999}. The present analytic question is narrower: regularity constraints select a small family of $\{\pm1\}$-valued chamber functions in the quasianalytic regime, while the non-quasianalytic regime permits them all.

\subsection{Established machinery versus the sign classification}
Stronger division, quasianalytic local-ring, resolution, and monomialization frameworks are available~\cite{Thilliez2008,BierstoneMilman2004,BelottoBierstoneChow2018,BelottoBierstone2023}. The proofs here deliberately avoid that machinery: the rigidity input is only the first nonzero Taylor polynomial, and the realization input is only the classical Denjoy--Carleman criterion together with cutoffs or positive-flat functions. This is the sense in which the note is self-contained without claiming foundational priority.

\section{Conclusion}

A first nonzero Taylor term provides a compact bridge from local regularity to global chamber signs. In finite order it yields explicit $C^d$ orthant and $C^N$ arrangement obstructions; in the smooth setting the same mechanism gives an elementary proof of the established signature-flatness consequence for affine arrangements. Combined with the Denjoy--Carleman criterion, the argument gives an exact realization classification: for $m$ distinct hyperplanes with $r$ chambers, the full normalized Roumieu class realizes exactly $2^{m+1}$ signed parity laws in the quasianalytic case and all $2^r$ chamber labelings in the non-quasianalytic case.

With fixed derivative bounds, the same vanishing jets control how rapidly an exact non-affine score approaches zero at the corner. Excluding wall strips changes the polynomial problem: at fixed dimension non-affine sign laws require logarithmic degree in the inverse width, whereas a fixed margin under a whole-cube amplitude bound requires reciprocal-width degree for every nonconstant law. The complete reductions make these distinctions explicit. The finite-order formulation, arrangement proof, and synthesis remain the contribution; the smooth obstruction, $R$-function constructions, positive-flat theory, and approximation inputs retain their classical attribution.

\begin{thebibliography}{99}
\small\raggedright

\bibitem{AcquistapaceBroglia1992}
F. Acquistapace and F. Broglia,
\newblock Signatures and flatness,
\newblock \emph{Journal f\"ur die reine und angewandte Mathematik} 425 (1992), 61--70.

\bibitem{AcquistapaceBroglia1993Correction}
F. Acquistapace and F. Broglia,
\newblock Correction to the paper ``Signatures and flatness'' (Vol. 425, pp. 61--70, 1992),
\newblock \emph{Journal f\"ur die reine und angewandte Mathematik} 444 (1993), 155--156.

\bibitem{AcquistapaceBroglia1994}
F. Acquistapace and F. Broglia,
\newblock More about signatures and approximation,
\newblock \emph{Geometriae Dedicata} 50 (1994), no. 2, 107--116.
\newblock DOI: \href{https://doi.org/10.1007/BF01265306}{10.1007/BF01265306}.

\bibitem{AcquistapaceBrogliaFortuna1994}
F. Acquistapace, F. Broglia, and E. Fortuna,
\newblock When is a distribution of signs locally completable?,
\newblock \emph{Canadian Journal of Mathematics} 46 (1994), no. 3, 449--473.
\newblock DOI: \href{https://doi.org/10.4153/CJM-1994-024-2}{10.4153/CJM-1994-024-2}.

\bibitem{Fortuna1991}
E. Fortuna,
\newblock Distributions de signes,
\newblock \emph{Mathematika} 38 (1991), no. 1, 50--62.
\newblock DOI: \href{https://doi.org/10.1112/S0025579300006434}{10.1112/S0025579300006434}.

\bibitem{BjornerEtAl1999}
A. Bj\"orner, M. Las Vergnas, B. Sturmfels, N. White, and G. M. Ziegler,
\newblock \emph{Oriented Matroids}, 2nd ed.,
\newblock Encyclopedia of Mathematics and its Applications, vol.~46, Cambridge University Press, Cambridge, 1999.

\bibitem{VarchenkoGelfand1987}
A. N. Varchenko and I. M. Gel'fand,
\newblock Heaviside functions of a configuration of hyperplanes,
\newblock \emph{Functional Analysis and Its Applications} 21 (1987), no. 4, 255--270.
\newblock DOI: \href{https://doi.org/10.1007/BF01077800}{10.1007/BF01077800}.

\bibitem{Moseley2017}
D. Moseley,
\newblock Equivariant cohomology and the Varchenko--Gelfand filtration,
\newblock \emph{Journal of Algebra} 472 (2017), 95--114.
\newblock DOI: \href{https://doi.org/10.1016/j.jalgebra.2016.10.010}{10.1016/j.jalgebra.2016.10.010}; arXiv:1110.5369.

\bibitem{BelottoBierstoneChow2018}
A. Belotto da Silva, E. Bierstone, and M. Chow,
\newblock Composite quasianalytic functions,
\newblock \emph{Compositio Mathematica} 154 (2018), no. 9, 1960--1973.
\newblock DOI: \href{https://doi.org/10.1112/S0010437X18007339}{10.1112/S0010437X18007339}; arXiv:1709.06629.

\bibitem{BelottoBierstone2023}
A. Belotto da Silva and E. Bierstone,
\newblock Monomialization of a quasianalytic morphism,
\newblock \emph{Annales Scientifiques de l'Ecole Normale Superieure} 56 (2023), no. 5, 1583--1651.
\newblock DOI: \href{https://doi.org/10.24033/asens.2562}{10.24033/asens.2562}; arXiv:1907.09502.

\bibitem{BierstoneMilman2004}
E. Bierstone and P. D. Milman,
\newblock Resolution of singularities in Denjoy--Carleman classes,
\newblock \emph{Selecta Mathematica} 10 (2004), no. 1, 1--28.
\newblock DOI: \href{https://doi.org/10.1007/s00029-004-0327-0}{10.1007/s00029-004-0327-0}; arXiv:math/0108204.


\bibitem{BrogliaTognoli1989}
F. Broglia and A. Tognoli,
\newblock Approximation of $C^\infty$-functions without changing their zero-set,
\newblock \emph{Annales de l'Institut Fourier} 39 (1989), no. 3, 611--632.
\newblock DOI: \href{https://doi.org/10.5802/aif.1178}{10.5802/aif.1178}.

\bibitem{Langenbruch2003}
M. Langenbruch,
\newblock A general approximation theorem of Whitney type,
\newblock \emph{Rev. R. Acad. Cien. Serie A. Mat. (RACSAM)} 97 (2003), no. 2, 287--303.
\newblock \url{https://rac.es/ficheros/doc/00128.pdf}.

\bibitem{Rainer2022}
A. Rainer,
\newblock Ultradifferentiable extension theorems: a survey,
\newblock \emph{Expositiones Mathematicae} 40 (2022), no. 3, 679--757.
\newblock DOI: \href{https://doi.org/10.1016/j.exmath.2021.12.001}{10.1016/j.exmath.2021.12.001}; arXiv:2107.01061.

\bibitem{Thilliez2008}
V. Thilliez,
\newblock On quasianalytic local rings,
\newblock \emph{Expositiones Mathematicae} 26 (2008), no. 1, 1--23.
\newblock DOI: \href{https://doi.org/10.1016/j.exmath.2007.04.001}{10.1016/j.exmath.2007.04.001}; arXiv:math/0509273.

\bibitem{Shapiro2007}
V. Shapiro,
\newblock Semi-analytic geometry with $R$-functions,
\newblock \emph{Acta Numerica} 16 (2007), 239--303.
\newblock DOI: \href{https://doi.org/10.1017/S096249290631001X}{10.1017/S096249290631001X}.

\bibitem{Sherstov2013}
A. A. Sherstov,
\newblock The intersection of two halfspaces has high threshold degree,
\newblock \emph{SIAM Journal on Computing} 42 (2013), no. 6, 2329--2374.
\newblock DOI: \href{https://doi.org/10.1137/100785260}{10.1137/100785260}.

\bibitem{KopotunPrymakShevchuk2018}
K. A. Kopotun, A. V. Prymak, and I. O. Shevchuk,
\newblock \emph{Introduction to Approximation Theory}, study guide, 24 January 2018,
\newblock Theorems 21.1 and 24.1.
\newblock \url{https://mathanalysis.knu.ua/wp-content/uploads/2020/07/at-study-guide-january-22-2018-1-1.pdf}.

\end{thebibliography}

\end{document}
